\chapter{工程部署与研究局限}
\section{集中式控制流程}
部署时，各 UAV 在时隙开始上报候选用户、任务元数据和估计 CSI；控制器构造初始化、扫描邻域、计算 deadline/min-CPU 准入，随后下发实际服务 UAV、子载波、角色、卸载比例和 CPU。若初始化求解、结构验证或硬 SIC 检查失败，控制器应回退到已验证的可行调度，而不是发布部分结果。

\section{协调与计算开销}
全局候选评价需要跨 UAV CSI 和同频占用信息。随着用户数增大，候选扫描时间明显上升；$N=50$ 的离线扩展性实验用于暴露这一趋势。可部署版本可先用干扰上界筛除候选，再对少量候选执行完整评价，但筛选规则必须保持主目标可审计。

\section{不完美信息}
正式鲁棒性套件分别改变 CSI 误差、SIC 残差、噪声和门限。规划 CSI 与真实 CSI 分离，最终 SIC 失败记为 outage，失败传输仍保留在其他链路干扰中。该模型未覆盖高速信道老化、同步误差和控制报文丢失。

\section{外部效度}
本文使用参数化 Monte Carlo 仿真，没有声称现场部署验证。公开车辆轨迹与地面行人终端场景不一致，因此未为了增加数据集而强行接入。后续可使用带步行标签的实测轨迹替代随机游走，但仍需解释多用户组合、服务区平移和任务生成假设。

\section{未优化变量}
无人机轨迹、发射功率、电池、长期任务队列、用户优先级和跨 UAV 任务迁移均未联合优化。深度强化学习代码属于历史探索范围，不作为本文核心基线。上述限制避免把一个时隙级组合调度器包装为完整自主空中网络。
